\originalchapter{连续介质与流体}
\originalsection{弹簧链、连续极限与场变分}
连续场不是把质点坐标随意换成函数. 原课从一条离散弹性链建立这种联系. 相邻质点平衡间隔为$a$, 第$i$个点的纵向位移为$\eta_i(t)$. 设线密度为$\lambda$, 每点质量$m=\lambda a$, 弹簧劲度$k=C/a$, 其中$C$为轴向刚度（弹性杆中可取Young模量乘截面积）. 离散拉格朗日量和方程为
\[
L_a=\sum_i\left[\frac{\lambda a}{2}\dot\eta_i^2
-\frac{C}{2a}(\eta_{i+1}-\eta_i)^2\right],\qquad
\lambda\ddot\eta_i=C\frac{\eta_{i+1}-2\eta_i+\eta_{i-1}}{a^2}.
\]
当位移变化的波长远大于$a$, 写$\eta_i(t)=\eta(x_i,t)$, 差商趋于$\eta_x$, 二阶差分趋于$\eta_{xx}$, $a\sum_i$趋于$\int\dd x$. 保持$\lambda,C$固定, 得
\[
L=\int\mathscr L\dd x,\qquad
\mathscr L=\frac\lambda2\eta_t^2-\frac C2\eta_x^2,
\qquad \lambda\eta_{tt}-C\eta_{xx}=0.
\]
波速$c=\sqrt{C/\lambda}$. 这里连续极限描述长波近似, 不保证有限$a$的链对任意短波也有相同色散. 例如离散平面波$\eta_i\propto e^{\mathrm i(kia-\omega t)}$给$\omega^2=4C\sin^2(ka/2)/(\lambda a^2)$, 仅在$|ka|\ll1$时接近$c^2k^2$.

同一场方程也可由驻定作用量得到. 对$S=\int\!\int\mathscr L(\eta,\eta_t,\eta_x,x,t)\dd x\dd t$, 在时间端点和固定空间边界取$\delta\eta=0$, 分别对$t,x$分部积分:
\[
\delta S=\int\!\int
\left(\mathscr L_\eta-\partial_t\mathscr L_{\eta_t}
-\partial_x\mathscr L_{\eta_x}\right)\delta\eta\dd x\dd t.
\]
因此场Euler--Lagrange方程为$\partial_t\mathscr L_{\eta_t}+\partial_x\mathscr L_{\eta_x}-\mathscr L_\eta=0$. 弹性链的密度给$\mathscr L_{\eta_t}=\lambda\eta_t$, $\mathscr L_{\eta_x}=-C\eta_x$, 立即重得波动方程. 原来每点一个方程, 现在每个空间标签$x$仍有一个方程, 空间导数表示邻点耦合. 三维时把$\partial_x\mathscr L_{\eta_x}$换成各空间方向之和; 多分量场则对每个分量独立变分. 这就是从有限自由度力学进入经典场论的具体入口.

\originalsection{流体场与质量守恒}
连续介质用场变量代替有限个质点坐标. 密度$\rho(\vect x,t)$给单位体积质量, 速度$\vect v(\vect x,t)$给经过该位置的微小流体团速度. 固定空间区域描述称Euler描述, 跟随流体团描述称Lagrange描述. 场$f$随流体团的变化是物质导数$Df/Dt=\partial_tf+\vect v\cdot\nabla f$, 这是把$f(\vect x(t),t)$用链式法则求导所得.

取固定体积$V$, 质量流出率为$\oint_{\partial V}\rho\vect v\cdot\vect n\dd A$, 因而$\dd\int_V\rho\dd V/\dd t=-\oint\rho\vect v\cdot\vect n\dd A$. 用散度定理, 任意体积给连续性方程$\partial_t\rho+\nabla\cdot(\rho\vect v)=0$. 若密度沿每个流体团不变, 即$D\rho/Dt=0$, 则$\nabla\cdot\vect v=0$. 不可压缩并不必然等于处处同一密度; 本册的定密度例子额外取$\rho$为常量.

理想流体无剪切应力, 面力为$-p\vect n$. 牛顿第二定律与连续性方程给Euler方程
\[
\rho\frac{D\vect v}{Dt}=-\nabla p+\rho\vect f,
\]
其中$\vect f$为单位质量体力. 若$\vect f=-\nabla\Phi$, 定常、不可压缩、定密度且沿同一流线, 将方程点乘$\dd\vect x\parallel\vect v$并积分得Bernoulli关系$p/\rho+v^2/2+\Phi=C$. 若进一步无旋, 用$(\vect v\cdot\nabla)\vect v=\nabla(v^2/2)-\vect v\times(\nabla\times\vect v)$可知常数在同一连通流区相同. 有旋流的不同流线一般不能共用同一$C$.

\originalsection{熵、动量通量与能量通量}
密度和速度给四个方程, 可压缩流体还有压强这个未知量. 理想流体无黏性、无热传导时, 每个流体团的比熵$s$满足$Ds/Dt=0$, 以状态方程$p=p(\rho,s)$闭合系统. 这里$s$是每单位质量的熵, 不把整团总熵与随质量标签不变的比熵混用.

连续性与Euler方程合并后, 动量第$i$分量满足
\[
\partial_t(\rho v_i)+\partial_j\Pi_{ij}=\rho f_i,
\qquad \Pi_{ij}=\rho v_iv_j+p\delta_{ij}.
\]
证明是把$\partial_t(\rho v_i)$按乘积展开, 用连续性消去$\partial_t\rho$, 再把$\rho Dv_i/Dt$替成$-\partial_ip+\rho f_i$. $\Pi_{ij}$表示第$i$方向动量通过法向第$j$方向的面的通量, 所以动量密度是向量而通量是张量. 其中$\rho v_iv_j$是流体搬运动量, $p\delta_{ij}$是压力传递动量. 它与只描述面力的Cauchy应力$-p\delta_{ij}$对象不同, 符号差异由其定义决定.

为写能量守恒, 引入比内能$e(\rho,s)$, 热力学微分为$\dd e=T\dd s+(p/\rho^2)\dd\rho$. 沿理想流体团$Ds/Dt=0$, $D\rho/Dt=-\rho\nabla\cdot\vect v$, 故$\rho De/Dt=-p\nabla\cdot\vect v$. 连续性将它变成
\[
\partial_t(\rho e)+\nabla\cdot(\rho e\vect v)
=-p\nabla\cdot\vect v.
\]
Euler方程点乘$\vect v$给动能变化中的$-\vect v\cdot\nabla p$. 两个压力项相加为$-\nabla\cdot(p\vect v)$: 压力既改变体团内能, 也通过体团边界作功. 若体力为$\vect f=-\nabla\Phi$, 总能量密度$w=\rho(v^2/2+e+\Phi)$满足
\[
\partial_tw+\nabla\cdot\bigl[(w+p)\vect v\bigr]
=\rho\partial_t\Phi.
\]
固定外势时右端为零, 积分后得到能量变化等于穿过边界的能量通量. $p/\rho$并非通常意义的比内能; 能量通量中的$e+p/\rho$称比焓, 这一项体现推动流体进入或离开区域所需的功.

若流动等熵且$p$可写成$\rho$的单值函数, 定义$h(\rho)=e+p/\rho$, 则$\dd h=\dd p/\rho$. 定常Euler方程沿流线积分给$v^2/2+h+\Phi=B$, 每条流线有自己的$B$. 常密度不可压缩模型中用$h=p/\rho+$常数, 重得前面的Bernoulli式. 如果有热交换、黏性损失或不同流线熵值, 就须按对应能量方程重新分析, 不能任意两点直接相减.

\originalsection{涡量与流体团的局部转动}
涡量为$\vect\Omega=\nabla\times\vect v$. Stokes积分定理给$\int_S\vect\Omega\cdot\vect n\dd A=\oint_{\partial S}\vect v\cdot\dd\vect l$, 将局部场与环量联系起来. 在一点展开速度梯度, 对称部分改变微小流体团的形状, 反对称部分产生刚体式转动, 局部角速度为$\vect\Omega/2$. 因此剪切流即使没有整体绕轴旋转, 也可有非零涡量.

取无旋体力$\vect f=-\nabla\Phi$. 对Euler方程取旋度, 再展开$\nabla\times(\vect v\times\vect\Omega)$, 可得
\[
\frac{D\vect\Omega}{Dt}
=(\vect\Omega\cdot\nabla)\vect v
-\vect\Omega(\nabla\cdot\vect v)
+\frac{\nabla\rho\times\nabla p}{\rho^2}.
\]
最后一项来自$-\nabla\times(\nabla p/\rho)$, 在等熵正压流体$p=p(\rho)$中为零. 常密度不可压缩理想流体仍有$D\vect\Omega/Dt=(\vect\Omega\cdot\nabla)\vect v$, 即涡线拉伸可以改变涡量. 原稿第123--124页把物质导数与旋度直接交换, 因而称涡量矢量沿流线恒定; 这一论证须修正. 正确的可用结论是: 对受无旋体力的平滑正压理想流体, 初始涡量为零时, 上述齐次输运方程保持它为零. 有黏性且$\rho,\nu$常数时另加$\nu\nabla^2\vect\Omega$, 会使涡量扩散.

\originalsection{静水与定常流的实际应用}
静水$v=0$给$\nabla p=-\rho\nabla\Phi$. 在重力$\Phi=gz$及常密度条件下$p(z)=p(0)-\rho gz$. 对浸没物体积分压力面力, 用散度定理得$\vect F_b=-\int_{V_{\rm displaced}}\nabla p\dd V=\rho gV_{\rm displaced}\hat{\vect z}$, 即Archimedes浮力. 若以等温理想气体$p=\rho R'T$为近似, 则$\dd p/\dd z=-gp/(R'T)$, $p(z)=p(0)e^{-gz/(R'T)}$. 这是所给状态方程与温度模型下的结果, 不应把任何温度梯度都称作必然不稳定.

定常流在各点不随时间变化, 不表示每个流体团都不加速. 变截面水管若近似不可压缩且截面速度均匀, 连续性给$A_1v_1=A_2v_2$, 同高度Bernoulli式给$p_1-p_2=\rho(v_2^2-v_1^2)/2$. 窄处速度较高、静压较低. 大水箱小孔出流时自由液面速度可忽略, 两点同为大气压, 故$v_{\rm jet}=\sqrt{2g(z_{\mathrm{surface}}-z_{\mathrm{hole}})}$. 忽略黏性和孔口收缩才有此理想值.

流线是在同一时刻处处沿$\vect v$的曲线; 迹线是同一个体团满足$\dot{\vect x}=\vect v(\vect x,t)$的运动轨迹. 定常时可消去$t$, 使二者重合; 非定常时一般不同. 驻点$v=0$处不能用$\dd y/\dd x=v_y/v_x$直接作除法. 在撞向驻点的流线上, $p_{\mathrm{stag}}-p_\infty=\rho v_\infty^2/2$. Pitot管比较驻点压与静压, 若压差以密度$\rho_l$的液柱在密度$\rho_a$的气流中测得, 则$\rho_av_\infty^2/2=(\rho_l-\rho_a)g\Delta h$. 液柱是压差读数, 并非直接等于速度.

\originalsection{位势流与圆柱}
不可压缩无旋流可局部写$\vect v=\nabla\phi$, 由连续性得$\nabla^2\phi=0$. 二维圆柱外的定常无环量流取远处$\vect v=U\hat{\vect x}$, 壁面法向速度为零. 与远处同样角依赖的Laplace解为$\phi=U(r+A/r)\cos\theta$. 在$r=R$令$\partial_r\phi=0$得$A=R^2$. 于是
\[
v_r=U(1-R^2/r^2)\cos\theta,\qquad
v_\theta=-U(1+R^2/r^2)\sin\theta.
\]
壁面速度为$-2U\sin\theta\hat{\vect\theta}$, Bernoulli式给$p(R,\theta)=p_\infty+\rho U^2(1-4\sin^2\theta)/2$. 压力对圆柱前后对称, 合力阻力为零. 这不是现实流动没有阻力: 黏性、边界层分离和尾流都被理想位势模型排除了. 此模型也假定零环量; 多连通圆柱外区域可另加环量流, 其单值势条件不同.

流函数为$\psi=U(r-R^2/r)\sin\theta=Uy(1-R^2/(x^2+y^2))$, 满足$v_r=r^{-1}\psi_\theta$, $v_\theta=-\psi_r$. 因$\vect v\cdot\nabla\psi=0$, $\psi$等值线就是流线. 壁面在$\psi=0$的一部分, 两个驻点在$x=\pm R,y=0$.

\originalsection{位势流的拓扑与理想球体流}
无旋流局部有速度势, 全域单值势还需检查区域的拓扑与环量. 原课几个基本解显示这一差别. 二维点源$\phi=A\log r$给$v_r=A/r$, 单位深度流量$2\pi A$; 三维点源$\phi=-A/r$给$v_r=A/r^2$, 总流量$4\pi A$. 它们只在$r>0$满足无源Laplace方程. 二维势涡$\phi=\Gamma\theta/(2\pi)$给$v_\theta=\Gamma/(2\pi r)$, 在原点外旋度为零, 环量却为$\Gamma$. 势随绕原点一圈增加$\Gamma$, 因而不是全域单值函数.

对常密度无旋流, Euler方程给$\nabla(\phi_t+v^2/2+p/\rho+\Phi)=0$, 故括号等于只依赖$t$的函数$b(t)$. 可以给$\phi$加纯时间函数来改变$b$, 速度仍不变; 对空间有变化的压强却不能用这一自由度消去. 这就是非定常位势流的Bernoulli关系.

原课的三维例子是半径$a$固定球体, 远处流体速度$U\hat{\vect z}$. 方程为$\nabla^2\phi=0$, 球面只要求不可穿透$\partial_r\phi|_a=0$, 远处$\phi\sim Ur\cos\theta$. 由于边界条件具有$\cos\theta$的角依赖, 选轴对称偶极解$\phi=U(r+A/r^2)\cos\theta$. 求径向导数并代入球面, 得$A=a^3/2$, 因而
\[
v_r=U(1-a^3/r^3)\cos\theta,\qquad
v_\theta=-U(1+a^3/(2r^3))\sin\theta.
\]
球面切向速度为$-3U\sin\theta/2$, 伯努利式给
\[
p(a,\theta)=p_\infty+\frac{\rho U^2}{8}(9\cos^2\theta-5).
\]
压强在$\theta$与$\pi-\theta$相同, 沿流向的压力积分为零. 这是d'Alembert零阻力悖论的球体实例. 现实中的无滑动条件不在这个理想模型中, 增大Re也不意味着可在壁面处随意去掉黏性.

\originalsection{声速、Mach数与扰动传播}
前述声波中的$c_s^2$可写为体积模量$B$除以平衡密度: $B=\rho_0(\partial p/\partial\rho)_s$, $c_s^2=B/\rho_0$. 线性Euler式还给$\partial_t(\nabla\times\vect v)=0$; 如果扰动初始无旋, 其速度同样满足波动方程. 压强、密度和纵向速度波由连续性及状态方程联系, 不是三个任意独立波形.

特征流速$U$与声速之比$\mathrm{Ma}=U/c_s$称Mach数. 低Mach、弱密度变化且无强热源的条件下, 声传播比流动变化快, 常能采用不可压缩近似. 单凭$\mathrm{Ma}\ll1$并不排除由温度或分层导致的密度差. 对匀速超声源, 时刻$t-\tau$发出的球形扰动半径为$c_s\tau$, 源走过距离为$U\tau$; 包络锥半角$\gamma$满足$\sin\gamma=c_s/U=1/\mathrm{Ma}$. 原课用这一几何图说明超声扰动不能传到上游. Mach锥的线性传播几何本身不是有限振幅冲击波跃变条件的证明.

\originalsection{声波}
取静止均匀平衡$\rho_0,p_0$, 令$\rho=\rho_0+\rho'$, $p=p_0+p'$, $\vect v$为一阶小量. 绝热局部状态关系为$p'=c_s^2\rho'$, 其中$c_s^2=(\partial p/\partial\rho)_s>0$. 连续性和Euler方程线性化为$\partial_t\rho'+\rho_0\nabla\cdot\vect v=0$, $\rho_0\partial_t\vect v=-\nabla p'$. 前一式对时间再微分并代入后一式, 得$\partial_t^2p'-c_s^2\nabla^2p'=0$. 声速是状态方程导数, 并非任意加上的传播常数.

对球对称标量, $\nabla^2p'=r^{-2}\partial_r(r^2\partial_rp')$. 令$w=rp'$, 在$r>0$有$\nabla^2p'=w_{rr}/r$, 所以出射球面波$p'=f(r-c_st)/r$满足方程. $1/r$衰减来自波面面积增长; $r=0$源点本身不在这段无源解的定义域.

\originalsection{黏性与量纲}
牛顿流体的黏性应力可写$\sigma_{ij}=-p\delta_{ij}+\eta(\partial_iv_j+\partial_jv_i)+(\zeta-2\eta/3)(\nabla\cdot\vect v)\delta_{ij}$. 定常物性、不可压缩时应力散度化为$-\nabla p+\eta\nabla^2\vect v$, 因而得到Navier--Stokes方程$\rho D\vect v/Dt=-\nabla p+\eta\nabla^2\vect v+\rho\vect f$. 此处$\eta$为动力黏度, $\nu=\eta/\rho$为运动黏度; 两者单位不同.

\begin{example}
半径$R$的长圆管内定常充分发展流, 速度为$v_z(r)$, 压强沿$z$均匀下降, 壁面无滑动. 求流量.
\end{example}
\begin{solution}
对流项为零, $0=-\partial_zp+\eta r^{-1}(rv'_z)'$. 积分给$rv'_z=(\partial_zp)r^2/(2\eta)+C$; 轴心有限排除$C/r$项, 再由$v_z(R)=0$得$v_z=-(\partial_zp)(R^2-r^2)/(4\eta)$. 流量$Q=\int_0^R2\pi rv_z\dd r=-\pi R^4\partial_zp/(8\eta)$. 这是层流、充分发展、牛顿流体模型, 入口区和湍流不能直接套用.
\end{solution}
尺度化$\vect x=L\tilde{\vect x}$, $\vect v=U\tilde{\vect v}$后, 惯性项与黏性项比为$\mathrm{Re}=\rho UL/\eta=UL/\nu$. 小Re时惯性可忽略, 称Stokes近似; 大Re时黏性在大尺度内部可能较小, 却仍能在边界层内重要.
\originalsection{黏性、边界条件与耗散}
把黏性应力单独记为$\tau_{ij}$, 则Cauchy应力为$-p\delta_{ij}+\tau_{ij}$, 动量通量为$\rho v_iv_j+p\delta_{ij}-\tau_{ij}$. 对各向同性牛顿流体, 令$S_{ij}=(\partial_iv_j+\partial_jv_i)/2$, $d=\nabla\cdot\vect v$, 有$\tau_{ij}=2\eta(S_{ij}-d\delta_{ij}/3)+\zeta d\delta_{ij}$. 对称速度梯度而非反对称部分进入此式, 是因为纯匀速平移或刚体转动没有邻层形变, 不应产生这种摩擦. 系数$\eta,\zeta\ge0$使
\[
\tau_{ij}\partial_jv_i
=2\eta\sum_{ij}(S_{ij}-d\delta_{ij}/3)^2+\zeta d^2\ge0.
\]
这表示机械能变成热的速率. 对不可压缩常密度流体, 无体力时, 动能密度$k=\rho v^2/2$满足
\[
\partial_tk+\partial_j\bigl[(k+p)v_j-v_i\tau_{ij}\bigr]
=-\tau_{ij}\partial_jv_i.
\]
积分后左侧的边界项是机械能通量. 若边界速度为零或在无穷远足够快衰减, 总动能导数为$-2\eta\int S_{ij}S_{ij}\dd V\le0$. 开放流动不能删掉边界通量, 例如压力泵持续输入能量才维持管流. 无热传导时$\rho T Ds/Dt=\tau_{ij}\partial_jv_i$, 故理想流体的熵守恒已被黏性产生的热改变.

无体力的常密度不可压缩流还可从Navier--Stokes取散度, 得压强Poisson方程$\nabla^2p=-\rho(\partial_iv_j)(\partial_jv_i)$. 这里用了$\partial_iv_i=0$, 说明压强是满足不可压缩约束所需的场, 不能先任意给压强再独立给速度. 固体壁面在黏性模型中通常要求无滑动$\vect v=\vect v_{\mathrm{wall}}$, 同时固定法向与切向速度. 这是与理想流体仅要求法向不可穿透的实质差别.

\originalsection{板间流与相似流动}
两平板在$y=0,h$, 定常充分发展速度为$v_x(y)$, 沿$x$压强梯度为$k<0$. 对流项$v_x\partial_xv_x$为零, Navier--Stokes给$\eta v_x''=k$. 两壁静止时$v_x(0)=v_x(h)=0$, 积分得$v_x=-ky(h-y)/(2\eta)$, 单位宽度流量为$-kh^3/(12\eta)$. 若$k=0$而上板以$U$运动, 条件为$v_x(0)=0$, $v_x(h)=U$, 得Couette流$v_x=Uy/h$. 一般二者可在这个线性约化模型中叠加. 对底壁的流向剪应力大小为$-kh/2$（静止两板）或$\eta U/h$（纯Couette流）; 必须明确是流体对壁还是壁对流体, 两者方向相反.

当几何只有一个尺度$L$、入流速度$U$且其他无量纲比也相同, 定常不可压缩流可写$\vect v/U=\vect h(\vect r/L,\mathrm{Re})$, $(p-p_\infty)/(\rho U^2)=g(\vect r/L,\mathrm{Re})$. 同Re的两个装置因而有相同无量纲场, 称动力相似. 这并不使Re成为任何几何下唯一参数: 若存在两个长度、可压缩性、自由表面或粗糙度, 其比值或其他无量纲数也须相同. 小Re问题的自然压力尺度改为$\eta U/L$, 它与惯性压力尺度相差$1/\mathrm{Re}$.

\originalsection{球体Stokes流的边值计算}
现取半径$a$固定球体, 无穷远$\vect v=\vect u=U\hat{\vect z}$, 球面无滑动$\vect v=0$. 常密度定常Navier--Stokes中惯性项尺度为$U^2/a$, 黏性项为$\nu U/a^2$, 比值为$\mathrm{Re}=Ua/\nu$. 当它很小, 主导边值问题为
\[
\nabla p=\eta\nabla^2\vect v,\qquad \nabla\cdot\vect v=0,
\qquad \vect v|_{r=a}=0,\quad\vect v\to\vect u\ (r\to\infty).
\]
散度给$\nabla^2p=0$. 边界问题对$\vect u$线性, 且只有球对称几何及入流方向, 取压强的衰减偶极形式$p=p_\infty+C\eta a(\vect u\cdot\vect r)/r^3$. 同样把速度写成
\[
\vect v=f(r)\vect u+k(r)(\vect u\cdot\vect r)\vect r.
\]
这两个矢量结构分别沿入流方向和径向; $f$无量纲, $k$量纲为长度的负二次方. 对该式求散度得$f'+r^2k'+4rk=0$. 再求Laplacian, 比较两个结构, 得
\[
\begin{aligned}
k''+\frac{6k'}r&=-\frac{3Ca}{r^5},\\
f''+\frac{2f'}r+2k&=\frac{Ca}{r^3}.
\end{aligned}
\]
第一式的解为$k=Ca/(2r^3)+A+B/r^5$. 远场要求$r^2k\to0$, 故$A=0$; 球面在所有角度为零要求$k(a)=f(a)=0$, 故$B=-Ca^3/2$. 将所得$k$代入散度式并以$f(\infty)=1$积分, 得
\[
k=\frac{Ca}{2r^3}\left(1-\frac{a^2}{r^2}\right),\qquad
f=1+\frac{Ca}{2r}+\frac{Ca^3}{6r^3}.
\]
最后$f(a)=0$给$C=-3/2$. 代回另一个ODE也成立, 故没有遗漏动量方程. 速度和压力为
\[
\begin{aligned}
v_r&=U\left(1-\frac{3a}{2r}+\frac{a^3}{2r^3}\right)\cos\theta,\\
v_\theta&=-U\left(1-\frac{3a}{4r}-\frac{a^3}{4r^3}\right)\sin\theta,\\
p&=p_\infty-\frac{3\eta Ua}{2r^2}\cos\theta.
\end{aligned}
\]
两速度分量在球面都为零, 远处恢复均匀流. 涡量为$\Omega_\phi=-3Ua\sin\theta/(2r^2)$, 并非理想位势流. 原稿第137页在$f$的积分常数和末次边界条件的文字上有排印矛盾, 上面保留$f(\infty)=1$, $f(a)=0$独立核对了全部条件.

\originalsection{压力与剪应力共同产生阻力}
球面外法向为$\hat{\vect r}$, 流体对球的牵引为$\vect t=(-pI+\tau)\hat{\vect r}$. 在球面,
\[
\tau_{rr}=2\eta\partial_rv_r=0,\qquad
\tau_{r\theta}=\eta\left(\frac1r\partial_\theta v_r+
\partial_rv_\theta-\frac{v_\theta}r\right)
=-\frac{3\eta U}{2a}\sin\theta.
\]
由于$\hat{\vect r}\cdot\hat{\vect z}=\cos\theta$, $\hat{\vect\theta}\cdot\hat{\vect z}=-\sin\theta$, 球受流向力为
\[
F_z=a^2\int_0^{2\pi}\!\int_0^\pi
[-p\cos\theta+\tau_{rr}\cos\theta-\tau_{r\theta}\sin\theta]
\sin\theta\dd\theta\dd\phi.
\]
均匀压强$p_\infty$贡献为零. 余下压力项为$2\pi\eta aU$, 切向黏性项为$4\pi\eta aU$, 总和为$6\pi\eta aU$. 因此静止流体中以$\vect V$运动的球受到
\[
\vect F_{\mathrm{drag}}=-6\pi\eta a\vect V.
\]
不能只积分压强而漏掉剪应力, 也不能只计算剪应力而忽略黏性流的非均匀压强. 稳态小Re解给阻力的主导项; 非定常附加质量、历史力、壁面效应与有限Re修正没有包含在该式中.

\begin{center}
\begin{tikzpicture}\begin{axis}[width=.46\textwidth,height=5cm,axis equal image,xmin=-4,xmax=4,ymin=-2.6,ymax=2.6,xtick={-4,-2,0,2,4},ytick={-2,0,2},xlabel={$z/a$},ylabel={$s/a$},title={理想位势流},tick label style={font=\scriptsize},title style={font=\small}]
\addplot[blue!65!black,no marks,thin] table[x=x,y=y,col sep=comma]{figures/sphere-ideal-0.csv};
\addplot[blue!65!black,no marks,thin] table[x=x,y expr=-\thisrow{y},col sep=comma]{figures/sphere-ideal-0.csv};
\addplot[blue!65!black,no marks,thin] table[x=x,y=y,col sep=comma]{figures/sphere-ideal-1.csv};
\addplot[blue!65!black,no marks,thin] table[x=x,y expr=-\thisrow{y},col sep=comma]{figures/sphere-ideal-1.csv};
\addplot[blue!65!black,no marks,thin] table[x=x,y=y,col sep=comma]{figures/sphere-ideal-2.csv};
\addplot[blue!65!black,no marks,thin] table[x=x,y expr=-\thisrow{y},col sep=comma]{figures/sphere-ideal-2.csv};
\addplot[blue!65!black,no marks,thin] table[x=x,y=y,col sep=comma]{figures/sphere-ideal-3.csv};
\addplot[blue!65!black,no marks,thin] table[x=x,y expr=-\thisrow{y},col sep=comma]{figures/sphere-ideal-3.csv};
\addplot[black,fill=gray!20,domain=0:360,samples=73] ({cos(x)},{sin(x)});
\draw[->,thick] (axis cs:-3.7,2.0)--(axis cs:-2.7,2.0) node[midway,above] {$\vect u$};
\end{axis}\end{tikzpicture}\hfill
\begin{tikzpicture}\begin{axis}[width=.46\textwidth,height=5cm,axis equal image,xmin=-4,xmax=4,ymin=-2.6,ymax=2.6,xtick={-4,-2,0,2,4},ytick={-2,0,2},xlabel={$z/a$},ylabel={$s/a$},title={Stokes黏性流},tick label style={font=\scriptsize},title style={font=\small}]
\addplot[blue!65!black,no marks,thin] table[x=x,y=y,col sep=comma]{figures/sphere-stokes-0.csv};
\addplot[blue!65!black,no marks,thin] table[x=x,y expr=-\thisrow{y},col sep=comma]{figures/sphere-stokes-0.csv};
\addplot[blue!65!black,no marks,thin] table[x=x,y=y,col sep=comma]{figures/sphere-stokes-1.csv};
\addplot[blue!65!black,no marks,thin] table[x=x,y expr=-\thisrow{y},col sep=comma]{figures/sphere-stokes-1.csv};
\addplot[blue!65!black,no marks,thin] table[x=x,y=y,col sep=comma]{figures/sphere-stokes-2.csv};
\addplot[blue!65!black,no marks,thin] table[x=x,y expr=-\thisrow{y},col sep=comma]{figures/sphere-stokes-2.csv};
\addplot[blue!65!black,no marks,thin] table[x=x,y=y,col sep=comma]{figures/sphere-stokes-3.csv};
\addplot[blue!65!black,no marks,thin] table[x=x,y expr=-\thisrow{y},col sep=comma]{figures/sphere-stokes-3.csv};
\addplot[black,fill=gray!20,domain=0:360,samples=73] ({cos(x)},{sin(x)});
\draw[->,thick] (axis cs:-3.7,2.0)--(axis cs:-2.7,2.0) node[midway,above] {$\vect u$};
\end{axis}\end{tikzpicture}
\end{center}

图的流线由轴对称流函数等值线实际求根得到. 理想球体流$\psi=Ur^2\sin^2\theta(1-a^3/r^3)/2$; Stokes流$\psi=Ur^2\sin^2\theta(1-3a/(2r)+a^3/(2r^3))/2$. 二者都不穿过球面, 但只有后者切向无滑动. 此图帮助读出绕流几何; 阻力结论来自上面的应力积分, 不能由流线形状直接猜得.

\originalsection{随Re改变的绕流与湍流概述}
原课最后讨论增大Re后的现象, 没有把完整非线性绕流解写成公式. 当$\mathrm{Re}\ll1$, 上述流线前后对称; Re增到约1时惯性修正使尾部开始不同; 数十到数百的范围可出现分离涡与回流, 进而失去定常性并周期脱落; 更大Re时尾流含许多尺度并呈不规则变化. 原稿按球半径的特征尺度举$10$--$100$、$150$--$1000$、$10^3$--$2\times10^5$等近似流态范围, 这些是所讨论装置的现象尺度, 不是由Re定义推出的普适临界数. 球面粗糙度、扰动、实验条件及用半径还是直径定义Re都会影响转变.

高Re时边界层转为湍流可能推迟分离, 使尾流变窄并降低总阻力; 因而“黏性耗散增强”与“某物体阻力总是增大”不是同一命题. 原课以球面凹坑和球类表面缝线为例说明这种作用. 管内湍流的瞬时速度不规则, 平均速度常比层流抛物线更平坦; 不能继续把充分发展层流公式当作其完整描述.

三维强湍流常有跨尺度能量传递: 大尺度输入的机械能传到更小涡, 最终由黏性变成热. 这是一种统计图景, 并不声称每个涡都按同一时序重复. 原课第7章进一步用量纲分析介绍惯性区间与Kolmogorov谱, 将在非线性章继续. 本册不以这些概述冒充湍流闭合理论或一般Navier--Stokes正则性证明.


本章来源: 8.09第6章第6.1--6.8节（原稿第104--141页）和作业8. 连续极限、守恒通量、理想球体流、黏性耗散、板间流与Stokes球体边值计算均据实际原稿中文重编; 圆柱例、流线计算及原稿涡量论证的修正为中文补充. 湍流在原课中是现象与标度入门, 没有完整闭合理论; 本册保留这一教学范围.
